\documentclass[11pt]{article}

\usepackage[utf8]{inputenc}
\usepackage[T1]{fontenc}
\usepackage{lmodern}
\usepackage{geometry}
\usepackage{amsmath,amssymb,amsfonts,amsthm,mathtools}
\usepackage{enumitem}
\usepackage{microtype}
\usepackage{hyperref}

\geometry{margin=1in}
\hypersetup{colorlinks=true, linkcolor=black, urlcolor=blue, citecolor=black}
\setlist[enumerate]{topsep=0.35em,itemsep=0.3em}
\setlist[itemize]{topsep=0.3em,itemsep=0.25em}

\newcommand{\Aether}{\AE{}ther}
\newcommand{\Aflow}{\AE{}ther-flow}
\newcommand{\TheoryProgram}{\Aether{} / \Aflow{} framework}
\newcommand{\Lang}{\mathcal{L}_{0}}
\newcommand{\LangPlus}{\mathcal{L}_{0}^{+H}}
\newcommand{\LedgerFam}{\mathfrak{H}}
\newcommand{\PiH}{\Pi_H}
\newcommand{\AdH}{\mathsf{Ad}_{H}}
\newcommand{\BlindH}{\mathsf{Blind}_{H}}
\newcommand{\GenH}{\mathsf{Gen}_{H}}
\newcommand{\EnvH}{\mathsf{Env}_{H}}
\newcommand{\EqSrc}{\mathsf{Eq}_{\mathrm{src}}}
\newcommand{\RetainH}{\mathsf{Retain}_{H}}
\newcommand{\NoTargetH}{\mathsf{NoTarget}_{H}}
\newcommand{\AuditH}{\mathsf{Audit}_{H}}
\newcommand{\Pass}{\mathsf{Pass}}
\newcommand{\Block}{\mathsf{Block}}
\newcommand{\Risk}{\mathsf{Risk}}

\newtheorem{auditfinding}{Audit finding}
\newtheorem{residualrisk}{Residual risk}
\newtheorem{separationtest}{Separation test}
\newtheorem{auditverdict}{Audit verdict}
\newtheorem{nextburden}{Next burden}

\title{\textbf{Localization Source-Basis Axiom Selector-Domain Ledger-Generation Ontology-Expansion Proposal Smuggling Audit}\\[0.35em]
\large Target-Import, Authority-Laundering, Quotient-Laundering, Retention, and Status Review}
\author{Smuggling Auditor AgentJob AJ-RT-20260613-015-001}
\date{June 13, 2026}

\begin{document}

\maketitle

\begin{abstract}
This draft/control audit reviews the noncanonical ledger-generation
ontology-expansion proposal for hidden target import, authority laundering,
quotient laundering, retention failure, and status laundering. The proposal
passes a narrow surface audit: it does not directly adopt target metric,
target proper-time, target observer protocol, target locality,
target-favorable rank, or benchmark-success evidence as a source primitive,
and it explicitly blocks canonical adoption. The pass is conditional and
non-promotional. The predicates \(\AdH\), \(\BlindH\), \(\EqSrc\),
\(\RetainH\), \(\NoTargetH\), and \(\AuditH\) remain theorem and
mechanization burdens, not discharged authority. Candidate reconstruction,
canonical ontology edit, benchmark promotion, Gate Chair review, and
completed-derivation language remain blocked.
\end{abstract}

\section{Scope}

This artifact is a draft/control output for task \texttt{RT-20260613-015}.
It audits
\texttt{30\_LOCALIZATION\_SOURCE\_BASIS\_AXIOM\_SELECTOR\_DOMAIN\_LEDGER\_GENERATION\_ONTOLOGY\_EXPANSION\_PROPOSAL.tex}.

The audit target is not whether the proposed extension should be adopted.
The target is narrower: whether the proposal, as a noncanonical object for
later scrutiny, imports target structure or silently expands protected
authority. The audited object is
\[
    \LangPlus =
    \Lang \cup
    \{\AdH,\BlindH,\GenH,\EqSrc,\RetainH,\NoTargetH,\AuditH\}.
\]

\section{Audit Method}

\begin{separationtest}[Five-channel smuggling screen]
The audit applies five checks:
\begin{enumerate}
    \item \textbf{Target import}: whether the proposal uses target metric,
    target proper-time, target observer protocol, target locality,
    target-favorable rank, or benchmark success as a source primitive.
    \item \textbf{Authority laundering}: whether a noncanonical proposal is
    treated as canonical ontology authority, Gate Chair authority, benchmark
    status, or candidate reconstruction permission.
    \item \textbf{Quotient laundering}: whether \(\EqSrc\) collapses adverse
    ledgers by definition before a source-equivalence theorem exists.
    \item \textbf{Retention failure}: whether \(\RetainH\) can omit,
    project away, or downgrade non-equivalent generated ledgers before a
    complete retention theorem exists.
    \item \textbf{Status laundering}: whether named predicates or audit
    markers are converted into theorem discharge, candidate support, or
    adoption authority.
\end{enumerate}
\end{separationtest}

\section{Target-Import Screen}

\begin{auditfinding}[No direct target primitive adoption]
The proposal does not define \(\GenH(\Pi_H)\) by target metric behavior,
target clock normalization, target observer protocols, target locality,
target-favorable ranks, or benchmark success. It names those channels only
inside \(\NoTargetH\), a proposed exclusion predicate.
\end{auditfinding}

\begin{residualrisk}[Negative target vocabulary must become source syntax]
The phrase ``exclude target evidence'' is not yet a mechanical source
definition. A later theorem or repair packet must convert \(\NoTargetH\)
into a source-syntax exclusion grammar or proof obligation. If \(\NoTargetH\)
is evaluated by knowing which ledger performs well against the target
benchmark, then the proposal fails. At this stage the predicate is an audit
target, not evidence.
\end{residualrisk}

\section{Authority-Laundering Screen}

\begin{auditfinding}[Canonical adoption is explicitly blocked]
The proposal states that \(\LangPlus\), \(\EnvH\), and \(\GenH\) are
noncanonical and require later audit plus human-gated ontology action before
any canonical use. It does not edit canonical ontology TeX and does not
request Gate Chair review.
\end{auditfinding}

\begin{residualrisk}[Audit marker is not authorization]
\(\AuditH\) correctly marks a required audit before adoption. It cannot be
used as a local source rule that authorizes adoption after this one audit.
Protected authority remains external to the artifact: canonical ontology
edit, benchmark promotion, Gate Chair verdict, and permanent authority
registration remain blocked unless a separate human-gated action authorizes
them.
\end{residualrisk}

\section{Quotient-Laundering Screen}

\begin{auditfinding}[Equivalence is placed after generation]
The proposal evaluates \(\EqSrc\) after generating the ledger family. That
ordering is correct: it prevents a definition from selecting a preferred
ledger before non-equivalent alternatives are visible.
\end{auditfinding}

\begin{residualrisk}[Equivalence semantics are underspecified]
\(\EqSrc(H',H'')\) is named as declared source relabeling, but the current
proposal does not supply the relabeling grammar, invariants, or failure
conditions. A quotient rule that erases adverse source behavior would
launder selection. Therefore any later use must prove that \(\EqSrc\) is an
equivalence relation and that each quotient class preserves the outputs and
statuses relevant to negative controls.
\end{residualrisk}

\section{Retention Screen}

\begin{auditfinding}[Family-first generation blocks immediate selection]
The proposal makes the first generated object a family
\(\LedgerFam_{\PiH}\), not a distinguished candidate-supporting ledger.
This is the correct local anti-smuggling shape because it avoids defining
\(H\) as the unique output by syntax.
\end{auditfinding}

\begin{residualrisk}[Retention is only a burden]
\(\RetainH\) is not yet a complete inventory theorem. The proposal states
that every non-equivalent generated ledger must be retained, but it does not
prove that \(\GenH(\Pi_H)\) is complete or that retained ledgers carry all
outputs and statuses needed for adverse comparison. Candidate support
therefore remains blocked.
\end{residualrisk}

\section{Status-Laundering Screen}

\begin{auditfinding}[The proposal labels theorem burdens as burdens]
The proposal keeps conservativity, completeness, equivalence, retention, and
target exclusion as theorem burdens. It does not present those burdens as
already discharged.
\end{auditfinding}

\begin{residualrisk}[Named predicates can be mistaken for solved structure]
The vocabulary \(\AdH\), \(\BlindH\), \(\GenH\), \(\EqSrc\), \(\RetainH\),
\(\NoTargetH\), and \(\AuditH\) is useful only if future work treats each
symbol as a formal obligation. If a later artifact treats the vocabulary
itself as proof of admissibility, blinding, equivalence, complete retention,
or target exclusion, it launders status.
\end{residualrisk}

\section{Audit Ledger}

\begin{center}
\begin{tabular}{p{0.2\linewidth}p{0.28\linewidth}p{0.28\linewidth}p{0.14\linewidth}}
\textbf{Channel} & \textbf{Local result} & \textbf{Residual risk} & \textbf{Status} \\
\hline
Target import & target predicates not adopted as source primitives & \(\NoTargetH\) needs mechanical source syntax & \(\Pass\), \(\Risk\) \\
Authority & canonical adoption explicitly blocked & \(\AuditH\) must not become permission & \(\Pass\), \(\Block\) \\
Quotient & equivalence applied after generation & \(\EqSrc\) may erase adverse behavior if underspecified & \(\Pass\), \(\Risk\) \\
Retention & family-first output blocks direct selection & completeness and output retention not proved & \(\Pass\), \(\Block\) \\
Status & theorem burdens remain labeled as burdens & named predicates may be mistaken for theorem evidence & \(\Pass\), \(\Block\) \\
\end{tabular}
\end{center}

\section{Local Verdict}

\begin{auditverdict}[Conditional local pass]
The ledger-generation ontology-expansion proposal passes the narrow
Smuggling Auditor screen as a noncanonical proposal. It is target-clean at
the level of declared intent and guard placement. It does not directly
adopt target structures as source primitives and does not claim protected
authority.
\end{auditverdict}

\begin{auditverdict}[No promotion or reconstruction follows]
The pass is not theorem evidence and not adoption authority. The proposal
still requires benchmark-independent stress testing of provenance packets,
admissibility rules, blinding records, source equivalence, retention
schemas, and target-exclusion grammars. Candidate reconstruction, canonical
ontology edit, benchmark promotion, Gate Chair review, and
completed-derivation language remain blocked.
\end{auditverdict}

\section{Next Burden}

\begin{nextburden}[Conservativity and benchmark-independence stress test]
The logical next role is a Refuter AgentJob that tests whether
\(\LangPlus\), \(\EnvH\), \(\GenH\), \(\EqSrc\), \(\RetainH\), and
\(\NoTargetH\) remain conservative and benchmark-independent under
admissible variation of provenance packets, admissibility rules, blinding
records, source-equivalence grammars, retention schemas, and target-exclusion
grammars. The test should ask whether favorable ledger selection reappears
under variation. It must not authorize candidate reconstruction, canonical
ontology edit, benchmark promotion, or Gate Chair review.
\end{nextburden}

\section*{References}

Ricciardi, A. S. (2026, June 13). \emph{Localization source-basis axiom
selector-domain ledger-generation ontology-expansion proposal: Noncanonical
source-language authority candidate for family-first ledger provenance}
[Unpublished manuscript]. The Aether-Flow Interpretation of Relativity
Research Project,
\texttt{research\_control/tasks/RT-20260613-010/artifacts/30\_LOCALIZATION\_SOURCE\_BASIS\_AXIOM\_SELECTOR\_DOMAIN\_LEDGER\_GENERATION\_ONTOLOGY\_EXPANSION\_PROPOSAL.tex}.

Ricciardi, A. S. (2026, June 13). \emph{Localization source-basis axiom
selector-domain ledger-provenance theorem attempt or controlled pause:
Single-surface current-language test for Pi\_H, H, and admissible ledger
retention} [Unpublished manuscript]. The Aether-Flow Interpretation of
Relativity Research Project,
\texttt{research\_control/tasks/RT-20260613-009/artifacts/29\_LOCALIZATION\_SOURCE\_BASIS\_AXIOM\_SELECTOR\_DOMAIN\_LEDGER\_PROVENANCE\_THEOREM\_ATTEMPT\_OR\_CONTROLLED\_PAUSE.tex}.

Ricciardi, A. S. (2026, June 13). \emph{Localization source-basis axiom
selector-domain source-forcing discharge protocol benchmark-independence
refutation: Variation test for four-cell discharge stability} [Unpublished
manuscript]. The Aether-Flow Interpretation of Relativity Research Project,
\texttt{research\_control/tasks/RT-20260613-008/artifacts/28\_LOCALIZATION\_SOURCE\_BASIS\_AXIOM\_SELECTOR\_DOMAIN\_SOURCE\_FORCING\_DISCHARGE\_PROTOCOL\_BENCHMARK\_INDEPENDENCE\_REFUTATION.tex}.

\end{document}
